\subsection{Các Use Case chi tiết}

Phần này đặc tả chi tiết từng chức năng đã trình bày ở biểu đồ Use Case tổng quát. Mỗi chức năng được mô tả bằng một biểu đồ Use Case chi tiết và một bảng đặc tả gồm tác nhân, điều kiện trước, điều kiện sau, luồng sự kiện chính và luồng sự kiện thay thế. Các đặc tả được trình bày theo thứ tự năm nhóm chức năng đã phân loại: nhóm dùng chung, nhóm quản lý người dùng, nhóm quản lý cuộc bầu cử, nhóm bỏ phiếu của cử tri và nhóm công khai.

\subsubsection{Xác thực người dùng}

Chức năng \textit{xác thực người dùng} là nền tảng cho mọi tương tác có liên quan tới quyền hạn trong hệ thống. Chức năng bao gồm ba thao tác gắn liền với vòng đời của một phiên làm việc: \textbf{đăng nhập} để khởi tạo phiên, \textbf{làm mới phiên} khi access token sắp hết hạn, và \textbf{đăng xuất} để chấm dứt phiên một cách chủ động.

\begin{figure}[H]
    \centering
    \makebox[\textwidth][c]{\includegraphics[width=0.91\textwidth, keepaspectratio]{Diagrams/Usecase/XacThucNguoiDung.pdf}}
    \caption{Biểu đồ Use Case ``Xác thực người dùng''}
    \label{fig:uc-xac-thuc-nguoi-dung}
\end{figure}

\begin{longtable}{|p{0.22\textwidth}|p{0.72\textwidth}|}
    \caption{Đặc tả Use Case ``Xác thực người dùng''}
    \label{tab:uc-xac-thuc-nguoi-dung} \\
    \hline
    \textbf{Thuộc tính} & \textbf{Nội dung} \\
    \hline
    \endfirsthead
    \hline
    \textbf{Thuộc tính} & \textbf{Nội dung} \\
    \hline
    \endhead
    \hline
    \endfoot
    \hline
    \endlastfoot
    Tên chức năng & Xác thực người dùng (đăng nhập, làm mới phiên, đăng xuất). \\
    \hline
    Tác nhân & Người dùng đã được cấp tài khoản trong hệ thống (quản trị viên hoặc cử tri). \\
    \hline
    Điều kiện trước &
    \begin{itemize}
        \item Tài khoản của người dùng đã được tạo và đang ở trạng thái kích hoạt.
        \item Đối với làm mới phiên và đăng xuất: người dùng đang sở hữu cặp token còn hiệu lực.
    \end{itemize} \\
    \hline
    Điều kiện sau &
    \begin{itemize}
        \item Đăng nhập thành công: người dùng nhận được cặp access token và refresh token mới, refresh token được lưu trong Redis.
        \item Làm mới phiên thành công: cặp token cũ bị thu hồi, cặp token mới được phát hành và ghi nhận.
        \item Đăng xuất thành công: refresh token được xóa khỏi Redis, phiên đăng nhập bị chấm dứt.
    \end{itemize} \\
    \hline
    Luồng sự kiện chính &
    \begin{enumerate}
        \item Người dùng gửi yêu cầu đăng nhập kèm thông tin tài khoản tới BFF Gateway.
        \item BFF Gateway chuyển tiếp yêu cầu tới Identity Service.
        \item Identity Service đối chiếu email và mật khẩu băm với cơ sở dữ liệu, sau đó phát hành cặp access token và refresh token; refresh token được lưu vào Redis gắn theo định danh người dùng.
        \item Người dùng nhận kết quả đăng nhập thành công.
        \item Khi access token sắp hết hạn, người dùng gửi refresh token để yêu cầu làm mới phiên; Identity Service đối chiếu refresh token lưu trong Redis, thu hồi token cũ và phát hành cặp token mới.
        \item Khi người dùng yêu cầu đăng xuất, Identity Service xóa refresh token tương ứng khỏi Redis và trả về kết quả thành công.
    \end{enumerate} \\
    \hline
    Luồng sự kiện thay thế &
    \begin{itemize}
        \item 3a. Sai thông tin đăng nhập: Identity Service trả về thông báo thông tin đăng nhập không hợp lệ và yêu cầu người dùng nhập lại.
        \item 3b. Tài khoản bị vô hiệu hóa: Identity Service từ chối cấp token và thông báo trạng thái tài khoản không khả dụng.
        \item 5a. Refresh token không khớp hoặc đã hết hạn: Identity Service từ chối làm mới phiên; người dùng phải đăng nhập lại từ đầu.
    \end{itemize} \\
    \hline
\end{longtable}

\subsubsection{Tạo tài khoản người dùng}

Chức năng \textit{tạo tài khoản người dùng} cho phép quản trị viên cấp tài khoản cho cử tri và ứng viên trước khi cuộc bầu cử diễn ra. Hệ thống hỗ trợ tạo đơn lẻ một tài khoản và tạo hàng loạt nhiều tài khoản, phù hợp cho các tổ chức cần nhập số lượng lớn cử tri.

\begin{figure}[H]
    \centering
    \makebox[\textwidth][c]{\includegraphics[width=0.7\textwidth, keepaspectratio]{Diagrams/Usecase/TaoDanhSachNguoiDung.pdf}}
    \caption{Biểu đồ Use Case ``Tạo tài khoản người dùng''}
    \label{fig:uc-tao-danh-sach-nguoi-dung}
\end{figure}

\begin{longtable}{|p{0.22\textwidth}|p{0.72\textwidth}|}
    \caption{Đặc tả Use Case ``Tạo tài khoản người dùng''}
    \label{tab:uc-tao-danh-sach-nguoi-dung} \\
    \hline
    \textbf{Thuộc tính} & \textbf{Nội dung} \\
    \hline
    \endfirsthead
    \hline
    \textbf{Thuộc tính} & \textbf{Nội dung} \\
    \hline
    \endhead
    \hline
    \endfoot
    \hline
    \endlastfoot
    Tên chức năng & Tạo tài khoản người dùng (đơn lẻ hoặc hàng loạt). \\
    \hline
    Tác nhân & Quản trị viên. \\
    \hline
    Điều kiện trước &
    \begin{itemize}
        \item Đã đăng nhập với vai trò quản trị viên.
        \item Danh sách email đầu vào không trùng với tài khoản hiện có trong hệ thống.
    \end{itemize} \\
    \hline
    Điều kiện sau & Một hoặc nhiều bản ghi tài khoản được tạo trong cơ sở dữ liệu Identity với mật khẩu đã được băm và trạng thái kích hoạt mặc định. \\
    \hline
    Luồng sự kiện chính &
    \begin{enumerate}
        \item Quản trị viên gửi yêu cầu tạo tài khoản đơn lẻ hoặc hàng loạt tới BFF Gateway kèm danh sách thông tin.
        \item BFF Gateway chuyển tiếp yêu cầu tới Identity Service.
        \item Identity Service kiểm tra tính hợp lệ của dữ liệu và ràng buộc trùng lặp email.
        \item Identity Service băm mật khẩu cho từng tài khoản và lưu hàng loạt vào cơ sở dữ liệu.
        \item Identity Service trả về danh sách bản ghi đã tạo (đơn lẻ) hoặc tổng số bản ghi đã tạo (hàng loạt).
        \item Quản trị viên nhận xác nhận tạo tài khoản thành công.
    \end{enumerate} \\
    \hline
    Luồng sự kiện thay thế &
    \begin{itemize}
        \item 1a. Dữ liệu đầu vào không hợp lệ: BFF Gateway từ chối yêu cầu, không gọi tới Identity Service.
        \item 3a. Có email trùng hoặc dữ liệu không hợp lệ: Identity Service tạo các bản ghi hợp lệ gửi về danh sách các bản ghi lỗi để quản trị viên chỉnh sửa.
    \end{itemize} \\
    \hline
\end{longtable}

\subsubsection{Tìm kiếm người dùng}

Chức năng \textit{tìm kiếm người dùng} hỗ trợ quản trị viên định vị tài khoản theo từ khóa, vai trò, trạng thái và phân trang, là cơ sở cho mọi thao tác quản lý người dùng tiếp theo.

\begin{figure}[H]
    \centering
    \makebox[\textwidth][c]{\includegraphics[width=0.7\textwidth, keepaspectratio]{Diagrams/Usecase/TimKiemNguoiDung.pdf}}
    \caption{Biểu đồ Use Case ``Tìm kiếm người dùng''}
    \label{fig:uc-tim-kiem-nguoi-dung}
\end{figure}

\begin{longtable}{|p{0.22\textwidth}|p{0.72\textwidth}|}
    \caption{Đặc tả Use Case ``Tìm kiếm người dùng''}
    \label{tab:uc-tim-kiem-nguoi-dung} \\
    \hline
    \textbf{Thuộc tính} & \textbf{Nội dung} \\
    \hline
    \endfirsthead
    \hline
    \textbf{Thuộc tính} & \textbf{Nội dung} \\
    \hline
    \endhead
    \hline
    \endfoot
    \hline
    \endlastfoot
    Tên chức năng & Tìm kiếm và lọc người dùng. \\
    \hline
    Tác nhân & Quản trị viên. \\
    \hline
    Điều kiện trước & Đã đăng nhập với vai trò quản trị viên. \\
    \hline
    Điều kiện sau & Hệ thống trả về danh sách người dùng phân trang theo các tiêu chí lọc. \\
    \hline
    Luồng sự kiện chính &
    \begin{enumerate}
        \item Quản trị viên gửi yêu cầu tìm kiếm tới BFF Gateway kèm các tham số (email, tên, vai trò, trạng thái, phân trang).
        \item BFF Gateway chuyển tiếp yêu cầu tới Identity Service.
        \item Identity Service xây dựng truy vấn lọc và phân trang trên cơ sở dữ liệu.
        \item Identity Service trả về tổng số bản ghi cùng trang dữ liệu tương ứng.
        \item Quản trị viên nhận kết quả tìm kiếm.
    \end{enumerate} \\
    \hline
    Luồng sự kiện thay thế &
    \begin{itemize}
        \item 1a. Dữ liệu đầu vào không hợp lệ: BFF Gateway từ chối yêu cầu, không gọi tới Identity Service.
    \end{itemize} \\
    \hline
\end{longtable}

\subsubsection{Quản lý trạng thái tài khoản}

Chức năng \textit{quản lý trạng thái tài khoản} cho phép quản trị viên kích hoạt hoặc vô hiệu hóa một tài khoản nhằm kiểm soát quyền truy cập của người dùng. Tài khoản đã bị vô hiệu hóa không thể đăng nhập và không thể tham gia bỏ phiếu kể cả khi đã thuộc danh sách cử tri của một cuộc bầu cử.

\begin{figure}[H]
    \centering
    \makebox[\textwidth][c]{\includegraphics[width=0.7\textwidth, keepaspectratio]{Diagrams/Usecase/QuanLyTrangThaiTaiKhoan.pdf}}
    \caption{Biểu đồ Use Case ``Quản lý trạng thái tài khoản''}
    \label{fig:uc-quan-ly-trang-thai-tai-khoan}
\end{figure}

\begin{longtable}{|p{0.22\textwidth}|p{0.72\textwidth}|}
    \caption{Đặc tả Use Case ``Quản lý trạng thái tài khoản''}
    \label{tab:uc-quan-ly-trang-thai-tai-khoan} \\
    \hline
    \textbf{Thuộc tính} & \textbf{Nội dung} \\
    \hline
    \endfirsthead
    \hline
    \textbf{Thuộc tính} & \textbf{Nội dung} \\
    \hline
    \endhead
    \hline
    \endfoot
    \hline
    \endlastfoot
    Tên chức năng & Kích hoạt hoặc vô hiệu hóa tài khoản người dùng. \\
    \hline
    Tác nhân & Quản trị viên. \\
    \hline
    Điều kiện trước &
    \begin{itemize}
        \item Đã đăng nhập với vai trò quản trị viên.
        \item Tài khoản cần thao tác tồn tại trong cơ sở dữ liệu.
    \end{itemize} \\
    \hline
    Điều kiện sau & Trạng thái của tài khoản được cập nhật; tài khoản nếu bị vô hiệu hóa không thể tiếp tục đăng nhập hoặc bỏ phiếu. \\
    \hline
    Luồng sự kiện chính &
    \begin{enumerate}
        \item Quản trị viên gửi yêu cầu kích hoạt hoặc vô hiệu hóa tới BFF Gateway theo định danh tài khoản.
        \item BFF Gateway chuyển tiếp yêu cầu tới Identity Service.
        \item Identity Service cập nhật trường trạng thái kích hoạt của bản ghi tài khoản.
        \item Identity Service trả về tài khoản sau cập nhật.
        \item Quản trị viên nhận xác nhận thao tác thành công.
    \end{enumerate} \\
    \hline
    Luồng sự kiện thay thế &
    \begin{itemize}
        \item 1a. Dữ liệu đầu vào không hợp lệ: BFF Gateway từ chối yêu cầu, không gọi tới Identity Service.
        \item 3a. Trạng thái yêu cầu trùng với trạng thái hiện tại: Identity Service trả về bản ghi không đổi, không phát sinh lỗi.
    \end{itemize} \\
    \hline
\end{longtable}

\subsubsection{Tạo cuộc bầu cử}

Chức năng \textit{tạo cuộc bầu cử} là điểm khởi đầu của vòng đời một cuộc bầu cử. Quản trị viên cần chỉ định tên duy nhất, danh sách ứng viên ban đầu và số ứng viên tối đa mỗi cử tri được chọn trên một lá phiếu (mặc định bằng $1$); cuộc bầu cử được khởi tạo ở trạng thái PENDING và sẽ chỉ trở nên khả dụng cho cử tri sau khi được điều hành chuyển sang trạng thái ACTIVE.

\begin{figure}[H]
    \centering
    \makebox[\textwidth][c]{\includegraphics[width=0.91\textwidth, keepaspectratio]{Diagrams/Usecase/TaoCuocBauCu.pdf}}
    \caption{Biểu đồ Use Case ``Tạo cuộc bầu cử''}
    \label{fig:uc-tao-cuoc-bau-cu}
\end{figure}

\begin{longtable}{|p{0.22\textwidth}|p{0.72\textwidth}|}
    \caption{Đặc tả Use Case ``Tạo cuộc bầu cử''}
    \label{tab:uc-tao-cuoc-bau-cu} \\
    \hline
    \textbf{Thuộc tính} & \textbf{Nội dung} \\
    \hline
    \endfirsthead
    \hline
    \textbf{Thuộc tính} & \textbf{Nội dung} \\
    \hline
    \endhead
    \hline
    \endfoot
    \hline
    \endlastfoot
    Tên chức năng & Khởi tạo một cuộc bầu cử mới. \\
    \hline
    Tác nhân & Quản trị viên. \\
    \hline
    Điều kiện trước &
    \begin{itemize}
        \item Đã đăng nhập với vai trò quản trị viên.
        \item Tên cuộc bầu cử dự kiến tạo chưa tồn tại trong hệ thống.
        \item Danh sách định danh ứng viên đầu vào đều thuộc vai trò Ứng viên và đang ở trạng thái kích hoạt.
        \item Số ứng viên tối đa được chọn nằm trong khoảng từ $1$ đến số lượng ứng viên trong danh sách.
    \end{itemize} \\
    \hline
    Điều kiện sau & Một cuộc bầu cử mới được tạo trong cơ sở dữ liệu Coordinator với trạng thái PENDING, danh sách ứng viên ban đầu và số ứng viên tối đa được chọn mỗi lá phiếu. \\
    \hline
    Luồng sự kiện chính &
    \begin{enumerate}
        \item Quản trị viên gửi yêu cầu tạo cuộc bầu cử tới BFF Gateway kèm tên, danh sách định danh ứng viên và số ứng viên tối đa được chọn mỗi lá phiếu.
        \item BFF Gateway chuyển tiếp yêu cầu tới Coordinator Service.
        \item Coordinator Service gửi yêu cầu xác minh danh sách ứng viên tới Identity Service.
        \item Identity Service trả về danh sách ứng viên hợp lệ tương ứng.
        \item Coordinator Service kiểm tra ràng buộc tên duy nhất và tính hợp lệ của số ứng viên tối đa được chọn, sau đó lưu bản ghi cuộc bầu cử với trạng thái PENDING.
        \item Quản trị viên nhận kết quả cuộc bầu cử đã được tạo.
    \end{enumerate} \\
    \hline
    Luồng sự kiện thay thế &
    \begin{itemize}
        \item 1a. Dữ liệu đầu vào không hợp lệ: BFF Gateway từ chối yêu cầu, không gọi tới Coordinator Service.
        \item 4a. Có định danh không thuộc vai trò Ứng viên hoặc đã bị vô hiệu hóa: Coordinator Service trả về số lượng định danh không hợp lệ.
        \item 5a. Tên cuộc bầu cử đã tồn tại: Coordinator Service trả về lỗi vi phạm ràng buộc duy nhất.
        \item 5b. Số ứng viên tối đa được chọn không hợp lệ: Coordinator Service từ chối tạo cuộc bầu cử và trả về thông báo lỗi.
    \end{itemize} \\
    \hline
\end{longtable}

\subsubsection{Tìm kiếm cuộc bầu cử}

Chức năng \textit{tìm kiếm cuộc bầu cử} cho phép quản trị viên định vị cuộc bầu cử theo từ khóa, trạng thái và khoảng thời gian. Đây là cơ sở chung cho các thao tác xem chi tiết, điều hành cũng như quản lý ứng viên và cử tri của cuộc bầu cử.

\begin{figure}[H]
    \centering
    \makebox[\textwidth][c]{\includegraphics[width=0.7\textwidth, keepaspectratio]{Diagrams/Usecase/TimKiemCuocBauCu.pdf}}
    \caption{Biểu đồ Use Case ``Tìm kiếm cuộc bầu cử''}
    \label{fig:uc-tim-kiem-cuoc-bau-cu}
\end{figure}

\begin{longtable}{|p{0.22\textwidth}|p{0.72\textwidth}|}
    \caption{Đặc tả Use Case ``Tìm kiếm cuộc bầu cử''}
    \label{tab:uc-tim-kiem-cuoc-bau-cu} \\
    \hline
    \textbf{Thuộc tính} & \textbf{Nội dung} \\
    \hline
    \endfirsthead
    \hline
    \textbf{Thuộc tính} & \textbf{Nội dung} \\
    \hline
    \endhead
    \hline
    \endfoot
    \hline
    \endlastfoot
    Tên chức năng & Tìm kiếm và lọc cuộc bầu cử. \\
    \hline
    Tác nhân & Quản trị viên. \\
    \hline
    Điều kiện trước & Đã đăng nhập với vai trò quản trị viên. \\
    \hline
    Điều kiện sau & Hệ thống trả về danh sách cuộc bầu cử phân trang theo các tiêu chí lọc. \\
    \hline
    Luồng sự kiện chính &
    \begin{enumerate}
        \item Quản trị viên gửi yêu cầu tìm kiếm tới BFF Gateway kèm các tham số (tên, trạng thái, khoảng thời gian, phân trang).
        \item BFF Gateway chuyển tiếp yêu cầu tới Coordinator Service.
        \item Coordinator Service xây dựng truy vấn lọc và phân trang trên cơ sở dữ liệu.
        \item Coordinator Service trả về tổng số bản ghi và trang dữ liệu tương ứng.
        \item Quản trị viên nhận kết quả tìm kiếm.
    \end{enumerate} \\
    \hline
    Luồng sự kiện thay thế &
    \begin{itemize}
        \item 1a. Dữ liệu đầu vào không hợp lệ: BFF Gateway từ chối yêu cầu, không gọi tới Coordinator Service.
    \end{itemize} \\
    \hline
\end{longtable}

\subsubsection{Xem chi tiết cuộc bầu cử}

Chức năng \textit{xem chi tiết cuộc bầu cử} trả về toàn bộ thông tin của một cuộc bầu cử, gồm các thuộc tính cơ bản (tên, trạng thái, mốc thời gian, mã giao dịch blockchain, gốc Merkle, public key tập thể), danh sách ứng viên và danh sách cử tri kèm hồ sơ chi tiết. Đây là chức năng nền tảng để quản trị viên kiểm tra cấu hình cuộc bầu cử trước khi điều hành.

\begin{figure}[H]
    \centering
    \makebox[\textwidth][c]{\includegraphics[width=0.91\textwidth, keepaspectratio]{Diagrams/Usecase/XemChiTietCuocBauCu.pdf}}
    \caption{Biểu đồ Use Case ``Xem chi tiết cuộc bầu cử''}
    \label{fig:uc-xem-chi-tiet-cuoc-bau-cu}
\end{figure}

\begin{longtable}{|p{0.22\textwidth}|p{0.72\textwidth}|}
    \caption{Đặc tả Use Case ``Xem chi tiết cuộc bầu cử''}
    \label{tab:uc-xem-chi-tiet-cuoc-bau-cu} \\
    \hline
    \textbf{Thuộc tính} & \textbf{Nội dung} \\
    \hline
    \endfirsthead
    \hline
    \textbf{Thuộc tính} & \textbf{Nội dung} \\
    \hline
    \endhead
    \hline
    \endfoot
    \hline
    \endlastfoot
    Tên chức năng & Xem chi tiết cuộc bầu cử kèm danh sách ứng viên và cử tri. \\
    \hline
    Tác nhân & Quản trị viên. \\
    \hline
    Điều kiện trước &
    \begin{itemize}
        \item Đã đăng nhập với vai trò quản trị viên.
        \item Cuộc bầu cử cần xem tồn tại trong cơ sở dữ liệu.
    \end{itemize} \\
    \hline
    Điều kiện sau & Hệ thống trả về thông tin đầy đủ của cuộc bầu cử cùng danh sách ứng viên và cử tri kèm hồ sơ. \\
    \hline
    Luồng sự kiện chính &
    \begin{enumerate}
        \item Quản trị viên gửi yêu cầu xem chi tiết tới BFF Gateway theo định danh cuộc bầu cử.
        \item BFF Gateway chuyển tiếp yêu cầu tới Coordinator Service.
        \item Coordinator Service truy xuất bản ghi cuộc bầu cử cùng danh sách định danh ứng viên và cử tri từ cơ sở dữ liệu.
        \item Coordinator Service yêu cầu Identity Service trả về hồ sơ chi tiết theo các định danh thu thập được.
        \item Coordinator Service tổng hợp thông tin và trả kết quả về BFF Gateway.
        \item Quản trị viên nhận chi tiết cuộc bầu cử cùng danh sách ứng viên và cử tri.
    \end{enumerate} \\
    \hline
    Luồng sự kiện thay thế &
    \begin{itemize}
        \item 1a. Định danh cuộc bầu cử không hợp lệ: BFF Gateway từ chối yêu cầu, không gọi tới Coordinator Service.
        \item 3a. Cuộc bầu cử không tồn tại: Coordinator Service trả về lỗi không tìm thấy.
    \end{itemize} \\
    \hline
\end{longtable}

\subsubsection{Quản lý ứng viên trong cuộc bầu cử}

Chức năng \textit{quản lý ứng viên trong cuộc bầu cử} cho phép quản trị viên thêm hoặc xóa ứng viên khỏi danh sách của một cuộc bầu cử cụ thể. Để tránh thay đổi cấu trúc cuộc bầu cử sau khi cử tri đã bắt đầu bỏ phiếu, chức năng này chỉ áp dụng cho cuộc bầu cử đang ở trạng thái PENDING.

\begin{figure}[H]
    \centering
    \makebox[\textwidth][c]{\includegraphics[width=0.91\textwidth, keepaspectratio]{Diagrams/Usecase/QuanLyUngVienTrongBauCu.pdf}}
    \caption{Biểu đồ Use Case ``Quản lý ứng viên trong cuộc bầu cử''}
    \label{fig:uc-quan-ly-ung-vien-trong-bau-cu}
\end{figure}

\begin{longtable}{|p{0.22\textwidth}|p{0.72\textwidth}|}
    \caption{Đặc tả Use Case ``Quản lý ứng viên trong cuộc bầu cử''}
    \label{tab:uc-quan-ly-ung-vien-trong-bau-cu} \\
    \hline
    \textbf{Thuộc tính} & \textbf{Nội dung} \\
    \hline
    \endfirsthead
    \hline
    \textbf{Thuộc tính} & \textbf{Nội dung} \\
    \hline
    \endhead
    \hline
    \endfoot
    \hline
    \endlastfoot
    Tên chức năng & Thêm hoặc xóa ứng viên trong cuộc bầu cử ở trạng thái PENDING. \\
    \hline
    Tác nhân & Quản trị viên. \\
    \hline
    Điều kiện trước &
    \begin{itemize}
        \item Đã đăng nhập với vai trò quản trị viên.
        \item Cuộc bầu cử ở trạng thái PENDING.
        \item Các định danh đầu vào đều thuộc vai trò Ứng viên và đang ở trạng thái ACTIVE.
    \end{itemize} \\
    \hline
    Điều kiện sau & Danh sách định danh ứng viên của cuộc bầu cử được cập nhật. \\
    \hline
    Luồng sự kiện chính &
    \begin{enumerate}
        \item Quản trị viên gửi yêu cầu thêm hoặc xóa ứng viên kèm danh sách định danh tới BFF Gateway.
        \item BFF Gateway chuyển tiếp yêu cầu tới Coordinator Service.
        \item Coordinator Service kiểm tra trạng thái cuộc bầu cử phải là PENDING.
        \item Coordinator Service xác minh danh sách định danh thông qua Identity Service.
        \item Coordinator Service xác nhận lại trạng thái PENDING của cuộc bầu cử rồi thêm hoặc xóa các bản ghi.
        \item Coordinator Service trả về cuộc bầu cử sau cập nhật.
        \item Quản trị viên nhận cuộc bầu cử sau cập nhật.
    \end{enumerate} \\
    \hline
    Luồng sự kiện thay thế &
    \begin{itemize}
        \item 1a. Dữ liệu đầu vào không hợp lệ: BFF Gateway từ chối yêu cầu, không gọi tới Coordinator Service.
        \item 3a. Cuộc bầu cử không ở trạng thái PENDING: Coordinator Service từ chối thao tác.
        \item 4a. Có định danh không hợp lệ trong danh sách đầu vào: Coordinator Service trả về số lượng định danh không hợp lệ.
    \end{itemize} \\
    \hline
\end{longtable}

\subsubsection{Quản lý cử tri trong cuộc bầu cử}

Chức năng \textit{quản lý cử tri trong cuộc bầu cử} cho phép quản trị viên cập nhật danh sách cử tri hợp lệ của một cuộc bầu cử. Vì danh sách cử tri quyết định ai có thể tham gia bỏ phiếu, hệ thống chỉ cho phép sửa đổi khi cuộc bầu cử còn ở trạng thái PENDING.

\begin{figure}[H]
    \centering
    \makebox[\textwidth][c]{\includegraphics[width=0.91\textwidth, keepaspectratio]{Diagrams/Usecase/QuanLyCuTriTrongBauCu.pdf}}
    \caption{Biểu đồ Use Case ``Quản lý cử tri trong cuộc bầu cử''}
    \label{fig:uc-quan-ly-cu-tri-trong-bau-cu}
\end{figure}

\begin{longtable}{|p{0.22\textwidth}|p{0.72\textwidth}|}
    \caption{Đặc tả Use Case ``Quản lý cử tri trong cuộc bầu cử''}
    \label{tab:uc-quan-ly-cu-tri-trong-bau-cu} \\
    \hline
    \textbf{Thuộc tính} & \textbf{Nội dung} \\
    \hline
    \endfirsthead
    \hline
    \textbf{Thuộc tính} & \textbf{Nội dung} \\
    \hline
    \endhead
    \hline
    \endfoot
    \hline
    \endlastfoot
    Tên chức năng & Thêm hoặc xóa cử tri trong cuộc bầu cử ở trạng thái PENDING. \\
    \hline
    Tác nhân & Quản trị viên. \\
    \hline
    Điều kiện trước &
    \begin{itemize}
        \item Đã đăng nhập với vai trò quản trị viên.
        \item Cuộc bầu cử ở trạng thái PENDING.
        \item Các định danh đầu vào đều thuộc vai trò cử tri và đang ở trạng thái hoạt động.
    \end{itemize} \\
    \hline
    Điều kiện sau & Danh sách định danh cử tri của cuộc bầu cử được cập nhật. \\
    \hline
    Luồng sự kiện chính &
    \begin{enumerate}
        \item Quản trị viên gửi yêu cầu thêm hoặc xóa cử tri kèm danh sách định danh tới BFF Gateway.
        \item BFF Gateway chuyển tiếp yêu cầu tới Coordinator Service.
        \item Coordinator Service xác minh danh sách định danh thông qua Identity Service.
        \item Coordinator Service xác nhận lại trạng thái PENDING của cuộc bầu cử rồi thêm hoặc xóa các bản ghi.
        \item Coordinator Service trả về cuộc bầu cử sau cập nhật.
        \item Quản trị viên nhận xác nhận thao tác thành công.
    \end{enumerate} \\
    \hline
    Luồng sự kiện thay thế &
    \begin{itemize}
        \item 1a. Dữ liệu đầu vào không hợp lệ: BFF Gateway từ chối yêu cầu, không gọi tới Coordinator Service.
        \item 3a. Có định danh không thuộc vai trò cử tri hoặc đã vô hiệu hóa: Coordinator Service trả về số lượng định danh không hợp lệ.
        \item 4a. Trạng thái cuộc bầu cử thay đổi trong khi giao dịch đang mở: Coordinator Service khôi phục giao dịch và trả về lỗi xung đột.
    \end{itemize} \\
    \hline
\end{longtable}

\subsubsection{Điều hành cuộc bầu cử}

Chức năng \textit{điều hành cuộc bầu cử} gắn liền với hai cột mốc chính trong vòng đời cuộc bầu cử: \textbf{bắt đầu} (chuyển sang trạng thái ACTIVE và sinh khóa công khai tập thể EC-Schnorr cho cơ chế chữ ký mù) và \textbf{đóng} (chuyển trạng thái sang CLOSED, xây dựng cây Merkle của toàn bộ cam kết phiếu và cam kết gốc Merkle lên blockchain). Cả hai thao tác đều cần phối hợp giữa nhiều service và mạng blockchain Hyperledger Fabric.

\begin{figure}[H]
    \centering
    \makebox[\textwidth][c]{\includegraphics[width=0.91\textwidth, keepaspectratio]{Diagrams/Usecase/DieuHanhCuocBauCu.pdf}}
    \caption{Biểu đồ Use Case ``Điều hành cuộc bầu cử''}
    \label{fig:uc-dieu-hanh-cuoc-bau-cu}
\end{figure}

\begin{longtable}{|p{0.22\textwidth}|p{0.72\textwidth}|}
    \caption{Đặc tả Use Case ``Điều hành cuộc bầu cử''}
    \label{tab:uc-dieu-hanh-cuoc-bau-cu} \\
    \hline
    \textbf{Thuộc tính} & \textbf{Nội dung} \\
    \hline
    \endfirsthead
    \hline
    \textbf{Thuộc tính} & \textbf{Nội dung} \\
    \hline
    \endhead
    \hline
    \endfoot
    \hline
    \endlastfoot
    Tên chức năng & Bắt đầu và đóng cuộc bầu cử. \\
    \hline
    Tác nhân & Quản trị viên (chính); Mạng blockchain (hệ thống ngoài). \\
    \hline
    Điều kiện trước &
    \begin{itemize}
        \item Đã đăng nhập với vai trò quản trị viên.
        \item Khi bắt đầu: cuộc bầu cử ở trạng thái PENDING, có tối thiểu hai ứng viên và ba cử tri.
        \item Khi đóng: cuộc bầu cử đang ở trạng thái ACTIVE.
    \end{itemize} \\
    \hline
    Điều kiện sau &
    \begin{itemize}
        \item Khi bắt đầu thành công: cuộc bầu cử chuyển sang trạng thái ACTIVE và lưu khóa công khai tập thể EC-Schnorr.
        \item Khi đóng thành công: cuộc bầu cử chuyển từ trạng thái ACTIVE sang CLOSING và cuối cùng là CLOSED; gốc Merkle và tham chiếu giao dịch blockchain được ghi nhận. Tự động phục hồi nếu tiến trình gặp sự cố ở trạng thái CLOSING.
    \end{itemize} \\
    \hline
    Luồng sự kiện chính &
    \begin{enumerate}
        \item Quản trị viên gửi yêu cầu bắt đầu hoặc đóng cuộc bầu cử tới BFF Gateway theo định danh.
        \item BFF Gateway chuyển tiếp yêu cầu tới Coordinator Service.
        \item (Bắt đầu) Coordinator Service kiểm tra trạng thái và số lượng tối thiểu của ứng viên, cử tri, sau đó gửi yêu cầu sinh cặp khóa song song tới ba Signing Node; ba node sinh khóa riêng, mã hóa và lưu vào cơ sở dữ liệu riêng của mỗi node, rồi trả về khóa công khai từng phần.
        \item (Bắt đầu) Coordinator Service cộng các khóa công khai để tạo khóa công khai tập thể, cập nhật trạng thái cuộc bầu cử sang ACTIVE và lưu khóa công khai vào cơ sở dữ liệu.
        \item (Đóng) Coordinator Service đọc toàn bộ cam kết phiếu đã xác nhận, xây dựng cây Merkle; ghi trạng thái trung gian CLOSING kèm gốc Merkle vào cơ sở dữ liệu trước, sau đó gửi yêu cầu cam kết gốc Merkle lên Mạng blockchain qua Chainlaunch.
        \item (Đóng) Chainlaunch ghi nhận gốc Merkle lên sổ cái và trả về mã giao dịch; Coordinator Service cập nhật trạng thái cuộc bầu cử từ CLOSING sang CLOSED kèm tham chiếu giao dịch blockchain được xác nhận.
        \item Quản trị viên nhận xác nhận thao tác thành công.
    \end{enumerate} \\
    \hline
    Luồng sự kiện thay thế &
    \begin{itemize}
        \item 3a. Không đủ số lượng ứng viên hoặc cử tri tối thiểu: Coordinator Service từ chối bắt đầu cuộc bầu cử.
        \item 3b. Một Signing Node phản hồi lỗi khi sinh khóa: Coordinator Service hủy thao tác và giữ cuộc bầu cử ở trạng thái PENDING.
        \item 5a. Chainlaunch phản hồi lỗi khi cam kết gốc Merkle: Coordinator Service truy vấn \pkg{GetMerkleRoot} trên blockchain để phân định. Nếu gốc Merkle đã được cam kết, hệ thống phục hồi trạng thái CLOSED kèm tham chiếu giao dịch và gốc Merkle từ blockchain. Nếu chưa cam kết, hệ thống khôi phục trạng thái ACTIVE và đặt lại giá trị của gốc Merkle.
    \end{itemize} \\
    \hline
\end{longtable}

\subsubsection{Xem danh sách phiếu trong cuộc bầu cử}

Chức năng \textit{xem danh sách phiếu} cho phép quản trị viên giám sát các phiếu đã nộp trong một cuộc bầu cử. Danh sách chỉ chứa các trường công khai như cam kết phiếu mù, tham chiếu giao dịch blockchain và thời điểm tạo; không có trường nào tiết lộ lựa chọn ứng viên của cử tri.

\begin{figure}[H]
    \centering
    \makebox[\textwidth][c]{\includegraphics[width=0.91\textwidth, keepaspectratio]{Diagrams/Usecase/XemDanhSachPhieu.pdf}}
    \caption{Biểu đồ Use Case ``Xem danh sách phiếu trong cuộc bầu cử''}
    \label{fig:uc-xem-danh-sach-phieu}
\end{figure}

\begin{longtable}{|p{0.22\textwidth}|p{0.72\textwidth}|}
    \caption{Đặc tả Use Case ``Xem danh sách phiếu trong cuộc bầu cử''}
    \label{tab:uc-xem-danh-sach-phieu} \\
    \hline
    \textbf{Thuộc tính} & \textbf{Nội dung} \\
    \hline
    \endfirsthead
    \hline
    \textbf{Thuộc tính} & \textbf{Nội dung} \\
    \hline
    \endhead
    \hline
    \endfoot
    \hline
    \endlastfoot
    Tên chức năng & Xem danh sách phiếu đã nộp trong một cuộc bầu cử. \\
    \hline
    Tác nhân & Quản trị viên. \\
    \hline
    Điều kiện trước &
    \begin{itemize}
        \item Đã đăng nhập với vai trò quản trị viên.
        \item Cuộc bầu cử tồn tại trong cơ sở dữ liệu.
    \end{itemize} \\
    \hline
    Điều kiện sau & Hệ thống trả về danh sách phiếu phân trang gồm cam kết phiếu, tham chiếu giao dịch blockchain và thời điểm tạo. \\
    \hline
    Luồng sự kiện chính &
    \begin{enumerate}
        \item Quản trị viên gửi yêu cầu xem danh sách phiếu tới BFF Gateway kèm tham số lọc và phân trang.
        \item BFF Gateway chuyển tiếp yêu cầu tới Coordinator Service.
        \item Coordinator Service kiểm tra cuộc bầu cử tồn tại, sau đó truy vấn danh sách phiếu theo tham số.
        \item Coordinator Service trả về danh sách phiếu cùng tổng số bản ghi.
        \item Quản trị viên nhận danh sách phiếu.
    \end{enumerate} \\
    \hline
    Luồng sự kiện thay thế &
    \begin{itemize}
        \item 1a. Dữ liệu đầu vào không hợp lệ: BFF Gateway từ chối yêu cầu, không gọi tới Coordinator Service.
        \item 3a. Cuộc bầu cử không tồn tại: Coordinator Service trả về lỗi không tìm thấy.
    \end{itemize} \\
    \hline
\end{longtable}

\subsubsection{Xem danh sách cuộc bầu cử của cử tri}

Chức năng \textit{xem danh sách cuộc bầu cử của cử tri} cho phép cử tri liệt kê các cuộc bầu cử mà mình được phép tham gia, cùng với trạng thái bỏ phiếu của bản thân trong từng cuộc bầu cử.

\begin{figure}[H]
    \centering
    \makebox[\textwidth][c]{\includegraphics[width=0.91\textwidth, keepaspectratio]{Diagrams/Usecase/XemBauCuCuaToi.pdf}}
    \caption{Biểu đồ Use Case ``Xem danh sách cuộc bầu cử của cử tri''}
    \label{fig:uc-xem-bau-cu-cua-toi}
\end{figure}

\begin{longtable}{|p{0.22\textwidth}|p{0.72\textwidth}|}
    \caption{Đặc tả Use Case ``Xem danh sách cuộc bầu cử của cử tri''}
    \label{tab:uc-xem-bau-cu-cua-toi} \\
    \hline
    \textbf{Thuộc tính} & \textbf{Nội dung} \\
    \hline
    \endfirsthead
    \hline
    \textbf{Thuộc tính} & \textbf{Nội dung} \\
    \hline
    \endhead
    \hline
    \endfoot
    \hline
    \endlastfoot
    Tên chức năng & Xem danh sách cuộc bầu cử mà cử tri tham gia. \\
    \hline
    Tác nhân & Cử tri. \\
    \hline
    Điều kiện trước & Đã đăng nhập với vai trò Cử tri. \\
    \hline
    Điều kiện sau & Hệ thống trả về danh sách cuộc bầu cử của chính cử tri kèm trạng thái bỏ phiếu của bản thân. \\
    \hline
    Luồng sự kiện chính &
    \begin{enumerate}
        \item Cử tri gửi yêu cầu xem danh sách cuộc bầu cử của mình tới BFF Gateway kèm tham số lọc tùy chọn theo trạng thái.
        \item BFF Gateway xác thực, lấy định danh cử tri từ phiên và chuyển tiếp yêu cầu tới Coordinator Service.
        \item Coordinator Service truy vấn các cuộc bầu cử mà cử tri được phép tham gia, kèm trạng thái bỏ phiếu của cử tri đó.
        \item Coordinator Service trả về danh sách cuộc bầu cử.
        \item Cử tri nhận kết quả.
    \end{enumerate} \\
    \hline
    Luồng sự kiện thay thế &
    \begin{itemize}
        \item 1a. Tham số lọc theo không hợp lệ: BFF Gateway từ chối yêu cầu, không gọi tới Coordinator Service.
    \end{itemize} \\
    \hline
\end{longtable}

\subsubsection{Xem chi tiết cuộc bầu cử đang tham gia}

Chức năng \textit{xem chi tiết cuộc bầu cử đang tham gia} dành cho cử tri muốn xem thông tin cụ thể của một cuộc bầu cử trước khi quyết định bỏ phiếu. Cử tri chỉ có thể xem chi tiết của các cuộc bầu cử mà mình thuộc danh sách cử tri hợp lệ.

\begin{figure}[H]
    \centering
    \makebox[\textwidth][c]{\includegraphics[width=0.91\textwidth, keepaspectratio]{Diagrams/Usecase/XemChiTietBauCuThamGia.pdf}}
    \caption{Biểu đồ Use Case ``Xem chi tiết cuộc bầu cử đang tham gia''}
    \label{fig:uc-xem-chi-tiet-bau-cu-tham-gia}
\end{figure}

\begin{longtable}{|p{0.22\textwidth}|p{0.72\textwidth}|}
    \caption{Đặc tả Use Case ``Xem chi tiết cuộc bầu cử đang tham gia''}
    \label{tab:uc-xem-chi-tiet-bau-cu-tham-gia} \\
    \hline
    \textbf{Thuộc tính} & \textbf{Nội dung} \\
    \hline
    \endfirsthead
    \hline
    \textbf{Thuộc tính} & \textbf{Nội dung} \\
    \hline
    \endhead
    \hline
    \endfoot
    \hline
    \endlastfoot
    Tên chức năng & Xem chi tiết một cuộc bầu cử của cử tri kèm trạng thái bỏ phiếu cá nhân. \\
    \hline
    Tác nhân & Cử tri. \\
    \hline
    Điều kiện trước &
    \begin{itemize}
        \item Đã đăng nhập với vai trò Cử tri.
        \item Cử tri thuộc danh sách cử tri của cuộc bầu cử cần xem.
    \end{itemize} \\
    \hline
    Điều kiện sau & Hệ thống trả về thông tin cuộc bầu cử, danh sách ứng viên và trạng thái bỏ phiếu của chính cử tri. \\
    \hline
    Luồng sự kiện chính &
    \begin{enumerate}
        \item Cử tri gửi yêu cầu xem chi tiết cuộc bầu cử tới BFF Gateway theo định danh.
        \item BFF Gateway xác thực, lấy định danh cử tri từ phiên và chuyển tiếp yêu cầu tới Coordinator Service.
        \item Coordinator Service truy xuất thông tin, danh sách định danh ứng viên của cuộc bầu cử mà cử tri yêu cầu.
        \item Coordinator Service yêu cầu Identity Service trả về hồ sơ ứng viên.
        \item Coordinator Service tổng hợp dữ liệu và trả về kết quả.
        \item Cử tri nhận chi tiết cuộc bầu cử cùng trạng thái bỏ phiếu của bản thân.
    \end{enumerate} \\
    \hline
    Luồng sự kiện thay thế &
    \begin{itemize}
        \item 3a. Cử tri không thuộc danh sách cử tri của cuộc bầu cử: Coordinator Service từ chối truy cập.
        \item 3b. Cuộc bầu cử không tồn tại: Coordinator Service trả về lỗi không tìm thấy.
    \end{itemize} \\
    \hline
\end{longtable}

\subsubsection{Bỏ phiếu}

Chức năng \textit{bỏ phiếu} là quy trình cốt lõi của hệ thống, được triển khai bằng cơ chế \textbf{chữ ký mù tập thể EC-Schnorr}. Cử tri lần lượt thực hiện ba bước: khởi tạo phiên bỏ phiếu để nhận cam kết tập thể từ ba Signing Node, gửi thông số ký mù để được ký từng phần và tổng hợp chữ ký, cuối cùng nộp cam kết phiếu mù lên blockchain. Sau ba bước, lá phiếu được ghi nhận trên blockchain mà không thành phần nào trong hệ thống biết được lựa chọn ứng viên của cử tri.

\begin{figure}[H]
    \centering
    \makebox[\textwidth][c]{\includegraphics[width=0.91\textwidth, keepaspectratio]{Diagrams/Usecase/BoPhieu.pdf}}
    \caption{Biểu đồ Use Case ``Bỏ phiếu''}
    \label{fig:uc-bo-phieu}
\end{figure}

\begin{longtable}{|p{0.22\textwidth}|p{0.72\textwidth}|}
    \caption{Đặc tả Use Case ``Bỏ phiếu''}
    \label{tab:uc-bo-phieu} \\
    \hline
    \textbf{Thuộc tính} & \textbf{Nội dung} \\
    \hline
    \endfirsthead
    \hline
    \textbf{Thuộc tính} & \textbf{Nội dung} \\
    \hline
    \endhead
    \hline
    \endfoot
    \hline
    \endlastfoot
    Tên chức năng & Bỏ phiếu mù theo cơ chế chữ ký mù tập thể EC-Schnorr. \\
    \hline
    Tác nhân & Cử tri (chính); Mạng blockchain (hệ thống ngoài). \\
    \hline
    Điều kiện trước &
    \begin{itemize}
        \item Đã đăng nhập với vai trò Cử tri.
        \item Cuộc bầu cử đang ở trạng thái hoạt động.
        \item Cử tri thuộc danh sách cử tri của cuộc bầu cử và chưa từng bỏ phiếu trong cuộc bầu cử này.
    \end{itemize} \\
    \hline
    Điều kiện sau &
    \begin{itemize}
        \item Cam kết phiếu mù của cử tri được ghi lên blockchain (trạng thái CONFIRMED) và lưu vào cơ sở dữ liệu kèm tham chiếu giao dịch; cử tri được đánh dấu đã bỏ phiếu và nhận về biên lai để xác minh sau này.
        \item Trong trường hợp lỗi tạm thời, tự động phục hồi bản ghi ở trạng thái {PENDING\_CHAIN} hoặc xóa bản ghi nếu phiếu chưa được ghi lên blockchain.
    \end{itemize} \\
    \hline
    Luồng sự kiện chính &
    \begin{enumerate}
        \item Cử tri gửi yêu cầu khởi tạo phiên bỏ phiếu tới BFF Gateway kèm định danh cuộc bầu cử.
        \item BFF Gateway xác thực, lấy định danh cử tri từ phiên và chuyển tiếp yêu cầu tới Coordinator Service.
        \item Coordinator Service kiểm tra điều kiện, gửi yêu cầu sinh cam kết song song tới ba Signing Node và tổng hợp thành cam kết tập thể; phiên được lưu trong Redis kèm thời gian sống ngắn.
        \item Cử tri nhận cam kết tập thể và khóa công khai tập thể, chọn từ một đến \pkg{maxSelectableCandidates} ứng viên, sau đó thực hiện làm mù phiếu (sinh các hệ số mù bí mật, tính cam kết mù và thông số ký mù gắn với tập ứng viên đã chọn).
        \item Cử tri gửi thông số ký mù tới BFF Gateway để yêu cầu ký phiếu mù.
        \item Coordinator Service kiểm tra phiên còn hợp lệ, gửi yêu cầu ký từng phần song song tới ba Signing Node; mỗi node ghi nhận một bản ghi đã ký vào cơ sở dữ liệu riêng và xóa nonce ngay sau khi ký.
        \item Coordinator Service tổng hợp các chữ ký từng phần thành chữ ký tổng và trả về cho cử tri; cử tri thực hiện giải mù để thu được chữ ký cuối cùng.
        \item Cử tri gửi cam kết phiếu mù cùng chữ ký tổng và định danh phiên tới BFF Gateway, và được chuyển tiếp yêu cầu tới Coordinator Service.
        \item Coordinator Service xác minh phiên và sinh định danh phiếu duy nhất, ghi trước bản ghi phiếu với trạng thái {PENDING\_CHAIN} vào cơ sở dữ liệu với ràng buộc duy nhất \pkg{(electionId, voterId)}, ngăn hai yêu cầu bỏ phiếu song song từ cùng một cử tri; sau đó gọi blockchain qua Chainlaunch để ghi cam kết phiếu.
        \item Chainlaunch ghi nhận phiếu vào sổ cái và trả về tham chiếu giao dịch; Coordinator Service cập nhật trạng thái phiếu từ {PENDING\_CHAIN} sang CONFIRMED kèm tham chiếu giao dịch và đánh dấu phiên đã hoàn tất.
        \item Cử tri nhận biên lai gồm định danh phiếu, cam kết phiếu mù và tham chiếu giao dịch blockchain.
    \end{enumerate} \\
    \hline
    Luồng sự kiện thay thế &
    \begin{itemize}
        \item 2a. Cử tri đã có bản ghi phiếu trong cuộc bầu cử: Coordinator Service từ chối khởi tạo phiên.
        \item 2b. Khóa công khai tập thể từ các Signing Node không khớp với khóa lưu trong cuộc bầu cử: Coordinator Service từ chối khởi tạo phiên nhằm chống tấn công lặp giữa các cuộc bầu cử.
        \item 5a. Một Signing Node phát hiện vi phạm ràng buộc duy nhất khi ghi bản ghi đã ký: yêu cầu ký toàn bộ bị từ chối.
        \item 8a. Bản ghi {PENDING\_CHAIN} cho cặp \pkg{(electionId, voterId)} đã tồn tại (yêu cầu bỏ phiếu song song hoặc cử tri đã bỏ phiếu): Coordinator Service từ chối với lỗi xung đột, không gọi blockchain.
        \item 8b. Chainlaunch phản hồi lỗi khi ghi phiếu: Coordinator Service truy vấn \texttt{GetVote} trên blockchain để phân định. Nếu phiếu đã được ghi, phục hồi trạng thái CONFIRMED kèm tham chiếu giao dịch từ blockchain. Nếu chưa ghi, xóa bản ghi {PENDING\_CHAIN} và trả về lỗi cho cử tri.
    \end{itemize} \\
    \hline
\end{longtable}

\subsubsection{Tiết lộ phiếu sau khi cuộc bầu cử đóng}

Chức năng \textit{tiết lộ phiếu} là bước cuối của quy trình bỏ phiếu mù. Sau khi cuộc bầu cử đóng, khách ẩn danh (do chính cử tri đại diện gửi mà không kèm phiên xác thực) công khai chữ ký EC-Schnorr đã giải mù cho tập ứng viên mình đã chọn thông qua Reveal-Vote Service không yêu cầu xác thực. Vì Reveal-Vote Service không biết người gửi là ai, điểm tin cậy duy nhất là tính hợp lệ của chữ ký EC-Schnorr đối với khóa công khai tập thể của cuộc bầu cử.

\begin{figure}[H]
    \centering
    \makebox[\textwidth][c]{\includegraphics[width=0.91\textwidth, keepaspectratio]{Diagrams/Usecase/TietLoPhieu.pdf}}
    \caption{Biểu đồ Use Case ``Tiết lộ phiếu sau khi cuộc bầu cử đóng''}
    \label{fig:uc-tiet-lo-phieu}
\end{figure}

\begin{longtable}{|p{0.22\textwidth}|p{0.72\textwidth}|}
    \caption{Đặc tả Use Case ``Tiết lộ phiếu sau khi cuộc bầu cử đóng''}
    \label{tab:uc-tiet-lo-phieu} \\
    \hline
    \textbf{Thuộc tính} & \textbf{Nội dung} \\
    \hline
    \endfirsthead
    \hline
    \textbf{Thuộc tính} & \textbf{Nội dung} \\
    \hline
    \endhead
    \hline
    \endfoot
    \hline
    \endlastfoot
    Tên chức năng & Tiết lộ phiếu sau khi cuộc bầu cử đóng. \\
    \hline
    Tác nhân & Khách (do chính cử tri đại diện gửi mà không kèm phiên xác thực); Mạng blockchain (hệ thống ngoài). \\
    \hline
    Điều kiện trước &
    \begin{itemize}
        \item Cuộc bầu cử ở trạng thái CLOSED.
        \item Khách sở hữu chữ ký EC-Schnorr đã giải mù cho lựa chọn ứng viên trong cuộc bầu cử.
    \end{itemize} \\
    \hline
    Điều kiện sau &
    \begin{itemize}
        \item Một bản ghi phiếu tiết lộ gồm định danh cuộc bầu cử, danh sách định danh ứng viên, khóa giải mù và tham số của chữ ký mù được lưu vào cơ sở dữ liệu Reveal-Vote và blockchain. Trong trường hợp blockchain đã ghi nhưng cơ sở dữ liệu chưa lưu, bản ghi được phục hồi qua truy vấn \texttt{GetUsedReveal}.
        \item Nếu tổng số phiếu đã tiết lộ đạt tổng số phiếu đã nộp, cuộc bầu cử được tự động chuyển sang trạng thái hoàn tất và các khóa của ba Signing Node được dọn dẹp.
    \end{itemize} \\
    \hline
    Luồng sự kiện chính &
    \begin{enumerate}
        \item Khách gửi yêu cầu tiết lộ phiếu tới Reveal-Vote Service qua cổng công khai, gồm danh sách định danh ứng viên đã chọn và chữ ký EC-Schnorr.
        \item Reveal-Vote Service yêu cầu Coordinator Service trả về thông tin cuộc bầu cử để xác nhận trạng thái CLOSED và lấy khóa công khai tập thể.
        \item Reveal-Vote Service chuẩn hóa (canonicalize) danh sách ứng viên, kiểm tra mọi ứng viên được chọn đều thuộc cuộc bầu cử và số lượng không vượt quá \pkg{maxSelectableCandidates}, sau đó xác minh chữ ký EC-Schnorr với khóa công khai tập thể và sinh khóa giải mù từ hai thành phần chữ ký.
        \item Reveal-Vote Service gửi yêu cầu ghi phiếu tiết lộ lên mạng blockchain qua Chainlaunch với revealKey chỉ được dùng một lần. Khi Chainlaunch phản hồi thành công, lưu bản ghi phiếu tiết lộ kèm tham chiếu giao dịch vào cơ sở dữ liệu.
        \item Reveal-Vote Service so sánh tổng số phiếu đã tiết lộ với tổng số phiếu đã nộp do Coordinator Service cung cấp, nếu bằng nhau thì Reveal-Vote Service yêu cầu Coordinator Service chuyển trạng thái cuộc bầu cử sang COMPLETED; Coordinator Service gửi thông điệp cho các Signing Node yêu cầu dọn dẹp khóa.
        \item Khách nhận xác nhận tiết lộ thành công kèm cờ cho biết cuộc bầu cử đã hoàn tất hay chưa.
    \end{enumerate} \\
    \hline
    Luồng sự kiện thay thế &
    \begin{itemize}
        \item 3a. Danh sách ứng viên không hợp lệ (vượt quá số tối đa cho phép hoặc chứa ứng viên không thuộc cuộc bầu cử): Reveal-Vote Service từ chối tiết lộ.
        \item 3b. Chữ ký EC-Schnorr không hợp lệ: Reveal-Vote Service từ chối tiết lộ.
        \item 4a. Khóa giải mù đã tồn tại trong cơ sở dữ liệu: vi phạm ràng buộc duy nhất \pkg{(electionId, revealKey)}, Reveal-Vote Service từ chối yêu cầu lặp nhằm chống tấn công phát lại (replay attack).
        \item 4b. Chainlaunch phản hồi lỗi khi ghi phiếu tiết lộ: Reveal-Vote Service truy vấn \texttt{GetUsedReveal} trên blockchain để phân định. Nếu revealKey đã tồn tại trên blockchain và \pkg{candidateIds} khớp, lưu bản ghi vào cơ sở dữ liệu với \texttt{blockchainRef}. Nếu revealKey chưa tồn tại trên blockchain, trả về lỗi cho khách.
    \end{itemize} \\
    \hline
\end{longtable}

\subsubsection{Sao lưu khóa bí mật phiếu}

Khóa bí mật phiếu (gồm các hệ số làm mù và hai thành phần chữ ký) được lưu cục bộ an toàn trên thiết bị của cử tri; việc mất thiết bị có thể dẫn đến mất khả năng tiết lộ và xác minh phiếu về sau nếu cử tri chưa sao lưu khóa. Chức năng \textit{sao lưu khóa bí mật phiếu} cho phép cử tri sao lưu các khóa này lên máy chủ theo mô hình \textit{zero-knowledge}: dữ liệu được mã hóa hoàn toàn phía thiết bị bằng khóa dẫn từ một mã PIN, máy chủ không bao giờ nắm được mã PIN hay nội dung khóa. Hình~\ref{fig:uc-quan-ly-khoa-bi-mat} mô tả hai chức năng sao lưu và khôi phục khóa, đều được cử tri thao tác sau khi đăng nhập.

\begin{figure}[H]
    \centering
    \makebox[\textwidth][c]{\includegraphics[width=0.7\textwidth, keepaspectratio]{Diagrams/Usecase/QuanLyKhoaBiMat.pdf}}
    \caption{Biểu đồ Use Case ``Sao lưu khóa bí mật phiếu''}
    \label{fig:uc-quan-ly-khoa-bi-mat}
\end{figure}

\begin{longtable}{|p{0.22\textwidth}|p{0.72\textwidth}|}
    \caption{Đặc tả Use Case ``Sao lưu khóa bí mật phiếu''}
    \label{tab:uc-sao-luu-khoa} \\
    \hline
    \textbf{Thuộc tính} & \textbf{Nội dung} \\
    \hline
    \endfirsthead
    \hline
    \textbf{Thuộc tính} & \textbf{Nội dung} \\
    \hline
    \endhead
    \hline
    \endfoot
    \hline
    \endlastfoot
    Tên chức năng & Sao lưu khóa bí mật phiếu lên máy chủ theo mô hình zero-knowledge. \\
    \hline
    Tác nhân & Cử tri. \\
    \hline
    Điều kiện trước &
    \begin{itemize}
        \item Đã đăng nhập với vai trò Cử tri.
        \item Thiết bị có ít nhất một khóa bí mật phiếu trong kho lưu trữ an toàn.
        \item Cử tri đặt một mã PIN dùng để dẫn xuất khóa mã hóa.
    \end{itemize} \\
    \hline
    Điều kiện sau & Một bản sao lưu đã mã hóa được lưu trên máy chủ (ghi đè bản cũ nếu có); máy chủ không nắm được mã PIN lẫn nội dung khóa. \\
    \hline
    Luồng sự kiện chính &
    \begin{enumerate}
        \item Cử tri chọn chức năng sao lưu khóa trên ứng dụng di động và nhập mã PIN.
        \item Ứng dụng đọc danh sách khóa bí mật phiếu từ kho lưu trữ an toàn cục bộ.
        \item Ứng dụng dẫn xuất khóa từ mã PIN bằng PBKDF2-SHA256 và mã hóa danh sách bằng AES-256-GCM, tạo thành một gói dữ liệu mã hóa (envelope).
        \item Ứng dụng gửi gói envelope kèm access token tới BFF Gateway; BFF Gateway xác thực rồi chuyển tiếp tới Identity Service.
        \item Identity Service áp dụng thêm một lớp mã hóa dữ liệu khi lưu trữ (\textit{at-rest encryption}) rồi ghi hoặc cập nhật bản sao lưu theo định danh cử tri.
        \item Cử tri nhận thông báo sao lưu thành công.
    \end{enumerate} \\
    \hline
    Luồng sự kiện thay thế &
    \begin{itemize}
        \item 2a. Thiết bị chưa có khóa bí mật nào: ứng dụng thông báo không có dữ liệu để sao lưu và dừng quy trình.
    \end{itemize} \\
    \hline
\end{longtable}

\subsubsection{Khôi phục khóa bí mật phiếu}

Chức năng \textit{khôi phục khóa bí mật phiếu} cho phép cử tri lấy lại các khóa đã sao lưu khi chuyển sang thiết bị mới hoặc cài lại ứng dụng. Quá trình giải mã diễn ra hoàn toàn phía thiết bị bằng mã PIN, do đó mã PIN sai sẽ làm bước xác thực của AES-256-GCM thất bại ngay tại máy khách.

\begin{longtable}{|p{0.22\textwidth}|p{0.72\textwidth}|}
    \caption{Đặc tả Use Case ``Khôi phục khóa bí mật phiếu''}
    \label{tab:uc-khoi-phuc-khoa} \\
    \hline
    \textbf{Thuộc tính} & \textbf{Nội dung} \\
    \hline
    \endfirsthead
    \hline
    \textbf{Thuộc tính} & \textbf{Nội dung} \\
    \hline
    \endhead
    \hline
    \endfoot
    \hline
    \endlastfoot
    Tên chức năng & Khôi phục khóa bí mật phiếu từ bản sao lưu trên máy chủ. \\
    \hline
    Tác nhân & Cử tri. \\
    \hline
    Điều kiện trước &
    \begin{itemize}
        \item Đã đăng nhập với vai trò Cử tri.
        \item Tồn tại một bản sao lưu khóa bí mật phiếu trên máy chủ.
        \item Cử tri nhập đúng mã PIN đã dùng khi sao lưu.
    \end{itemize} \\
    \hline
    Điều kiện sau & Khóa bí mật phiếu được giải mã và ghi vào kho lưu trữ an toàn của thiết bị. \\
    \hline
    Luồng sự kiện chính &
    \begin{enumerate}
        \item Cử tri chọn chức năng khôi phục khóa trên ứng dụng di động và nhập mã PIN.
        \item Ứng dụng gửi yêu cầu tải bản sao lưu kèm access token; BFF Gateway xác thực rồi chuyển tiếp tới Identity Service.
        \item Identity Service truy vấn bản sao lưu theo định danh cử tri, gỡ lớp mã hóa dữ liệu khi lưu trữ và trả về gói envelope do máy khách tạo.
        \item Ứng dụng dẫn xuất khóa từ mã PIN và giải mã gói envelope bằng AES-256-GCM.
        \item Ứng dụng ghi khóa bí mật phiếu vào kho lưu trữ an toàn và thông báo khôi phục thành công.
    \end{enumerate} \\
    \hline
    Luồng sự kiện thay thế &
    \begin{itemize}
        \item 3a. Chưa có bản sao lưu trên máy chủ: Identity Service thông báo không tìm thấy dữ liệu để khôi phục.
        \item 4a. Mã PIN không đúng: bước xác thực AES-256-GCM thất bại, ứng dụng thông báo ``Mã PIN không đúng''.
    \end{itemize} \\
    \hline
\end{longtable}

\subsubsection{Xác minh phiếu bầu}

Chức năng \textit{xác minh phiếu bầu} cho phép bất kỳ khách nào dùng biên lai để kiểm tra xem một phiếu cụ thể đã thực sự được ghi nhận vào cuộc bầu cử hay chưa. Quy trình xác minh diễn ra đa tầng: đối chiếu với cơ sở dữ liệu, đối chiếu với blockchain và (sau khi cuộc bầu cử đóng) kiểm tra bằng chứng Merkle giữa cả hai nguồn. Kết quả chỉ được coi là hợp lệ khi mọi tầng đều xác nhận khớp.

\begin{figure}[H]
    \centering
    \makebox[\textwidth][c]{\includegraphics[width=0.91\textwidth, keepaspectratio]{Diagrams/Usecase/XacMinhPhieuBau.pdf}}
    \caption{Biểu đồ Use Case ``Xác minh phiếu bầu''}
    \label{fig:uc-xac-minh-phieu-bau}
\end{figure}

\begin{longtable}{|p{0.22\textwidth}|p{0.72\textwidth}|}
    \caption{Đặc tả Use Case ``Xác minh phiếu bầu''}
    \label{tab:uc-xac-minh-phieu-bau} \\
    \hline
    \textbf{Thuộc tính} & \textbf{Nội dung} \\
    \hline
    \endfirsthead
    \hline
    \textbf{Thuộc tính} & \textbf{Nội dung} \\
    \hline
    \endhead
    \hline
    \endfoot
    \hline
    \endlastfoot
    Tên chức năng & Xác minh phiếu bầu đa tầng. \\
    \hline
    Tác nhân & Khách; Mạng blockchain (hệ thống ngoài). \\
    \hline
    Điều kiện trước &
    \begin{itemize}
        \item Cuộc bầu cử tồn tại; phiếu cần xác minh đã được nộp vào cuộc bầu cử.
        \item Khách sở hữu biên lai gồm định danh phiếu, cam kết phiếu mù và tham chiếu giao dịch blockchain.
    \end{itemize} \\
    \hline
    Điều kiện sau & Hệ thống trả về kết quả xác minh từng tầng kèm kết luận tổng hợp. \\
    \hline
    Luồng sự kiện chính &
    \begin{enumerate}
        \item Khách gửi yêu cầu xác minh phiếu tới BFF Gateway qua điểm cuối công khai kèm thông tin từ biên lai.
        \item BFF Gateway chuyển tiếp yêu cầu tới Coordinator Service.
        \item Coordinator Service kiểm tra cuộc bầu cử tồn tại, sau đó đối chiếu bản ghi phiếu trong cơ sở dữ liệu với biên lai và truy vấn bản ghi phiếu trên Mạng blockchain để kiểm tra tham chiếu giao dịch và cam kết.
        \item Nếu cuộc bầu cử đã đóng hoặc đã hoàn tất, Coordinator Service xây dựng bằng chứng Merkle từ tập cam kết phiếu, kiểm tra cục bộ tính hợp lệ của bằng chứng, đối chiếu gốc Merkle với gốc trên blockchain và yêu cầu blockchain xác minh bằng chứng.
        \item Coordinator Service tổng hợp kết quả xác minh ở từng tầng và trả về kết luận chỉ coi là hợp lệ khi mọi tầng đều khớp.
        \item Khách nhận kết quả xác minh chi tiết.
    \end{enumerate} \\
    \hline
    Luồng sự kiện thay thế &
    \begin{itemize}
        \item 3a. Bản ghi phiếu không tồn tại hoặc không khớp biên lai: trả về kết quả không hợp lệ tại tầng cơ sở dữ liệu.
        \item 4a. Cuộc bầu cử chưa đóng: bỏ qua xác minh Merkle, chỉ trả kết quả của hai tầng đầu.
        \item 4b. Gốc Merkle giữa cơ sở dữ liệu và blockchain không khớp: trả về kết quả không hợp lệ tại tầng Merkle.
    \end{itemize} \\
    \hline
\end{longtable}

\subsubsection{Xem kết quả bầu cử}

Chức năng \textit{xem kết quả bầu cử} công bố số phiếu đã được tính cho từng ứng viên sau khi cuộc bầu cử đóng. Vì hệ thống lưu kết quả đồng thời trong cơ sở dữ liệu và trên blockchain, kết quả trả về kèm cả hai con số để khách có thể đối chiếu chéo và phát hiện ngay nếu có sự sai khác.

\begin{figure}[H]
    \centering
    \makebox[\textwidth][c]{\includegraphics[width=0.91\textwidth, keepaspectratio]{Diagrams/Usecase/XemKetQuaBauCu.pdf}}
    \caption{Biểu đồ Use Case ``Xem kết quả bầu cử''}
    \label{fig:uc-xem-ket-qua-bau-cu}
\end{figure}

\begin{longtable}{|p{0.22\textwidth}|p{0.72\textwidth}|}
    \caption{Đặc tả Use Case ``Xem kết quả bầu cử''}
    \label{tab:uc-xem-ket-qua-bau-cu} \\
    \hline
    \textbf{Thuộc tính} & \textbf{Nội dung} \\
    \hline
    \endfirsthead
    \hline
    \textbf{Thuộc tính} & \textbf{Nội dung} \\
    \hline
    \endhead
    \hline
    \endfoot
    \hline
    \endlastfoot
    Tên chức năng & Xem kết quả tổng hợp số phiếu cho từng ứng viên. \\
    \hline
    Tác nhân & Khách; Mạng blockchain (hệ thống ngoài). \\
    \hline
    Điều kiện trước & Cuộc bầu cử ở trạng thái đóng hoặc hoàn tất. \\
    \hline
    Điều kiện sau & Hệ thống trả về số phiếu cho từng ứng viên kèm đối chiếu giữa cơ sở dữ liệu và blockchain. \\
    \hline
    Luồng sự kiện chính &
    \begin{enumerate}
        \item Khách gửi yêu cầu xem kết quả tới Reveal-Vote qua điểm cuối HTTP công khai theo định danh cuộc bầu cử.
        \item Reveal-Vote yêu cầu Coordinator Service trả về thông tin cuộc bầu cử để kiểm tra trạng thái.
        \item Reveal-Vote thực hiện song song ba thao tác: nhóm và đếm số phiếu đã tiết lộ theo ứng viên từ cơ sở dữ liệu; truy vấn kết quả kiểm phiếu từ Mạng blockchain; lấy hồ sơ ứng viên từ Identity Service.
        \item Reveal-Vote ghép tên ứng viên với hai con số đếm và đối chiếu tổng số phiếu giữa cơ sở dữ liệu và blockchain.
        \item Khách nhận kết quả chi tiết cho từng ứng viên.
    \end{enumerate} \\
    \hline
    Luồng sự kiện thay thế &
    \begin{itemize}
        \item 2a. Cuộc bầu cử chưa đóng: Reveal-Vote từ chối trả kết quả.
    \end{itemize} \\
    \hline
\end{longtable}

\subsubsection{Kiểm toán số phiếu}

Chức năng \textit{kiểm toán số phiếu} cho phép bất kỳ khách nào lấy về bản đối chiếu các con số quan trọng của một cuộc bầu cử giữa cơ sở dữ liệu nội bộ và blockchain. Kết quả gồm tổng số phiếu đã nộp, số phiếu đã tiết lộ và trạng thái cam kết gốc Merkle, nhờ đó mọi bên quan tâm có thể tự kiểm chứng tính nhất quán của hệ thống mà không cần truy cập sâu vào dữ liệu nghiệp vụ.

\begin{figure}[H]
    \centering
    \makebox[\textwidth][c]{\includegraphics[width=0.91\textwidth, keepaspectratio]{Diagrams/Usecase/KiemToanSoPhieu.pdf}}
    \caption{Biểu đồ Use Case ``Kiểm toán số phiếu''}
    \label{fig:uc-kiem-toan-so-phieu}
\end{figure}

\begin{longtable}{|p{0.22\textwidth}|p{0.72\textwidth}|}
    \caption{Đặc tả Use Case ``Kiểm toán số phiếu''}
    \label{tab:uc-kiem-toan-so-phieu} \\
    \hline
    \textbf{Thuộc tính} & \textbf{Nội dung} \\
    \hline
    \endfirsthead
    \hline
    \textbf{Thuộc tính} & \textbf{Nội dung} \\
    \hline
    \endhead
    \hline
    \endfoot
    \hline
    \endlastfoot
    Tên chức năng & Kiểm toán số phiếu giữa cơ sở dữ liệu và blockchain. \\
    \hline
    Tác nhân & Khách; Mạng blockchain (hệ thống ngoài). \\
    \hline
    Điều kiện trước & Cuộc bầu cử tồn tại trong hệ thống. \\
    \hline
    Điều kiện sau & Hệ thống trả về bản đối chiếu các con số liên quan đến cuộc bầu cử giữa cơ sở dữ liệu và blockchain. \\
    \hline
    Luồng sự kiện chính &
    \begin{enumerate}
        \item Khách gửi yêu cầu kiểm toán tới Reveal-Vote qua điểm cuối HTTP công khai theo định danh cuộc bầu cử.
        \item Reveal-Vote thực hiện song song ba thao tác: đếm số phiếu đã tiết lộ trong cơ sở dữ liệu của Reveal-Vote; yêu cầu Coordinator Service trả về tổng số phiếu đã nộp; truy vấn các chỉ số kiểm toán từ Mạng blockchain gồm tổng số phiếu, số phiếu đã tiết lộ và trạng thái cam kết gốc Merkle.
        \item Reveal-Vote tổng hợp các con số và trạng thái, sau đó trả kết quả đối chiếu cho khách.
        \item Khách nhận bảng đối chiếu chi tiết.
    \end{enumerate} \\
    \hline
    Luồng sự kiện thay thế &
    \begin{itemize}
        \item 2a. Cuộc bầu cử không tồn tại: Reveal-Vote trả về lỗi không tìm thấy.
    \end{itemize} \\
    \hline
\end{longtable}
